% TOP_PROBLEM: {"id": 30007458, "title": "FPRAS for perfect matchings in general graphs", "category_id": 15, "category": {"id": 15, "name": "computer_science", "display_name": "Computer Science", "description": "Computational complexity, algorithms, and theoretical CS.", "slug": "computer-science", "order_index": 15, "created_at": "2026-07-31T15:26:25.671Z"}, "status": "open"}
% FIRST_POSTED: 2026-10-06
% ENTRY_KIND: statement_only
% !TeX program = lualatex
\documentclass[11pt]{article}
\usepackage[margin=25mm]{geometry}
\usepackage{fontspec}
\setmainfont{Latin Modern Roman}
\usepackage{amsmath,amssymb,mathtools}
\usepackage{unicode-math}
\setmathfont{Latin Modern Math}
\usepackage{xurl,hyperref}
\hypersetup{colorlinks=true,urlcolor=blue}
\setcounter{secnumdepth}{0}
\setlength{\parindent}{0pt}
\setlength{\parskip}{5pt}
\setlength{\emergencystretch}{3em}
\newcommand{\R}{\mathbb R}
\newcommand{\N}{\mathbb N}
\newcommand{\Z}{\mathbb Z}
\newcommand{\Q}{\mathbb Q}
\newcommand{\C}{\mathbb C}
\newcommand{\D}{\mathbb D}
\newcommand{\F}{\mathbb F}
\newcommand{\sref}[2]{\hyperlink{source:#1:#2}{[S#2]}}
\newcommand{\cataloguescope}[1]{#1}
\begin{document}
\section*{FPRAS for perfect matchings in general graphs}\label{30007458}
\subsection{Definitions and mathematical statement}
\paragraph{Statement recorded in the supplied source.}
Can the number of perfect matchings of an arbitrary explicitly given finite general graph be estimated to relative error $\varepsilon $ with failure probability $\delta $ in worst-case bit time polynomial in input length, $\varepsilon ^{-1}$ and $\log (\delta ^{-1})$, with zero returned certainly when there are no perfect matchings?
\subsection{Short English statement}
Can the number of perfect matchings of an arbitrary explicitly given finite general graph be estimated to relative error $\varepsilon $ with failure probability $\delta $ in worst-case bit time polynomial in input length, $\varepsilon ^{-1}$ and $\log (\delta ^{-1})$, with zero returned certainly when there are no perfect matchings?
\subsection{Sources}
\begin{itemize}
\item[S1] OpenAI, manuscript citation page. Source of the formulation motivating this local problem.
\url{https://github.com/openai/math/blob/adc7f1241b42e322a6451854ab7e4b4c146bf78a/preprints/A-Fully-Polynomial-Randomized-Approximation-Scheme-for-Perfect-Matchings-in-General-Graphs-September-23-2026/README.md}.
\end{itemize}
\end{document}
